\section*{Chimie(a points) : Acide propanoique - Synthese d'un ester}
La racine de la canne à sucre est caractérisée par l'odeur d'un composé organique naturel ( $E$ ) appartenant à la famille des esters. Les chimistes ont pu synthétiser l'ester ( $E$ ) à partir de l'acidepropanoїque et de l'alcool (2-methylpropan-1-ol).\\
Cet exercice vise :

\begin{itemize}
  \item l'étude d'une solution aqueuse d'acide propanoїque ;
  \item l'étude de la synthèse de l'ester ( $E$ ).
\end{itemize}

\section*{Données :}
\begin{center}
\begin{tabular}{|l|c|c|c|}
\hline
Composé organique & \begin{tabular}{c}
Acide \\
propanoïque \\
\end{tabular} & 2-methylpropan-1-ol & Ester $(E)$ \\
\hline
Formule chimique & $\mathrm{C}_{2} \mathrm{H}_{5} \mathrm{CO}_{2} \mathrm{H}$ & $\mathrm{C}_{4} \mathrm{H}_{9} \mathrm{OH}$ &  \\
\hline
Masse molaire moléculaire &  &  & $130 \mathrm{~g} . \mathrm{mol}^{-1}$ \\
\hline
\end{tabular}
\end{center}

\section*{1. Étude d'une solution aqueuse d'acide propanoïque}
On considère une solution aqueuse ( $S_{1}$ ) de volume $V_{1}=1 L$ contenant la quantité de matière initiale $n_{1}=2.10^{-3} \mathrm{~mol}$ d'acide propanoïque. Le $p H$ de cette solution est $p H_{1}=3,78$ à $25^{\circ} \mathrm{C}$.\\
1.1. Écrire l'équation chimique modélisant la réaction de l'acide propanoïque avec l'eau.\\
1.2. Dresser le tableau d'avancement de la réaction en utilisant les grandeurs $n_{1}$ et $x_{\text {éq }}$ l'avancement de la réaction à l'état d'équilibre du système chimique.\\
1.3. Calculer la valeur de $x_{\dot{e} q}$.\\
1.4. Calculer la valeur de $\tau_{1}$, taux d'avancement final de la réaction. Conclure.\\
1.5. Montrer que l'expression de la constante d'acidité $K_{A}$ du couple ( $\mathrm{C}_{2} \mathrm{H}_{5} \mathrm{CO}_{2} \mathrm{H}_{\text {(aq) }} / \mathrm{C}_{2} \mathrm{H}_{5} \mathrm{CO}_{2(a q)}^{-}$) s'écrit : $K_{A}=\frac{x_{\dot{e} q}^{2}}{V_{1}\left(n_{1}-x_{\dot{e} q}\right)}$. Calculer sa valeur.\\
1.6. Représenter le diagramme de prédominance des espèces du couple ( $\mathrm{C}_{2} \mathrm{H}_{5} \mathrm{CO}_{2} \mathrm{H}_{(a q)} / \mathrm{C}_{2} \mathrm{H}_{5} \mathrm{CO}_{2(a q)}^{-}$).\\
1.7. Indiquer, en justifiant, l'espèce prédominante parmi $\mathrm{C}_{2} \mathrm{H}_{5} \mathrm{CO}_{2} \mathrm{H}_{(a q)}$ et $\mathrm{C}_{2} \mathrm{H}_{5} \mathrm{CO}_{2(a q)}^{-}$dans la solution étudiée ( $S_{1}$ ).\\
1.8. On considère une autre solution aqueuse ( $S_{2}$ ) d'acide propanoïque de concentration molaire $C_{2}=1.10^{-3} \mathrm{~mol} . L^{-1}$. Le taux d'avancement final $\tau_{2}$ de la réaction de l'acide propanoïque avec l'eau dans ce cas est $\tau_{2}=11,6 \%$. En comparant $\tau_{1}$ et $\tau_{2}$ conclure.

\section*{2. Étude de la synthèse de l'ester ( $E$ )}
Pour synthétiser l'ester $(E)$, on introduit dans un ballon les quantités de matière $n_{1}=0,6 \mathrm{~mol}$ d'acide propanoïque et $n_{2}=0,6 \mathrm{~mol}$ de l'alcool (2-methylpropan-1-ol), quelques gouttes d'acide sulfurique concentré et quelques grains de pierre ponce.\\
On chauffe à reflux le système chimique pendant une durée suffisante pour atteindre l'état d'équilibre du système. La masse finale de l'ester $(E)$ formé est $m(E)=52 g$.\\
2.1. Donner les caractéristiques de la réaction d'estérification.\\
2.2. Justifier le choix du chauffage à reflux.\\
2.3. Donner le rôle de l'acide sulfurique concentré au cours de cette synthèse.\\
2.4. Écrire la formule semi-développée de l'ester ( $E$ ).\\
2.5. Calculer la valeur du rendement $r$ de cette synthèse.\\
2.6. Kecopier sur votre copie le numero de la question et ecrıre la lettre correspondante a la seule proposition vraie :\\
Pour synthétiser l'ester ( $E$ ) selon une transformation rapide et totale, il faut :

\begin{center}
\begin{tabular}{|l|l|}
\hline
$\mathbf{A}$ & Augmenter la température du milieu réactionnel \\
\hline
$\mathbf{B}$ & Remplacer l'alcool par la solution d'hydroxyde de sodium $\left(\mathrm{Na}_{(a q)}^{+}+\mathrm{HO}_{(a q)}^{-}\right)$ \\
\hline
$\mathbf{C}$ & Remplacer l'acide propanoïque par l'anhydride propanoïque \\
\hline
$\mathbf{D}$ & Utiliser 1,2 mol de l'acide et 1,2 mol de l'alcool \\
\hline
\end{tabular}
\end{center}


